% !TeX root = ../../Book1.tex
% 纲要标题：### 14.4 函数域与超越次数
% 纲要位置：计划纲要.md，第 520 行。

\section{函数域与超越次数}
\label{b1:sec:1404}


% 纲要内容（仅作写作计划，尚非教材正文）：
%
% * 有理函数域与代数函数域
% * 次数与超越次数的不同角色；剩余有限次数随超越基的选择而变
% * 第三卷维数理论的域论背景
%
% ---
%

从多项式环取分式域，就得到可以进行除法的函数表达式。变量之间若有代数关系，取分式以后关系仍然存在。本节用几个可直接计算的例子说明：扩张次数描述代数方程的复杂度，超越次数则记录独立参数的数目。

\subsection{有理函数域与代数函数域}
\label{b1:subsec:1404:01}

\begin{definition}\label{b1:def:function-field}
设 \(L/k\) 为域扩张。如果 \(L/k\) 有限生成，则称 \(L\) 为 \(k\) 上的\term[hanshu-yu]{函数域}[function field]。本书允许函数域的超越次数为零，因而有限代数扩张也属于函数域。设 \(r\geq0\) 为整数；如果 \(t_1,\ldots,t_r\in L\) 在 \(k\) 上代数独立，则称 \(k(t_1,\ldots,t_r)\) 为 \(k\) 上的\term[youli-hanshu-yu]{有理函数域}[rational function field]；有理函数域的有限代数扩张通常称为\term[daishu-hanshu-yu]{代数函数域}[algebraic function field]。若 \(L\) 中所有在 \(k\) 上代数的元素都属于 \(k\)，则称 \(k\) 在 \(L\) 中相对代数闭。本定义不要求这一条件；它也不同于要求 \(k\) 自身是代数闭域。
\end{definition}
由\cref{b1:prop:finitely-generated-field-structure}，每个函数域都是某个有理函数域的有限扩张。例如 \(\mathbb Q(\sqrt2)/\mathbb Q\) 的超越次数为零，仍符合本书的函数域定义；此时超越基为空，参数域就是 \(\mathbb Q\)。

\ProseParagraphBreak
相对代数闭只比较 \(k\) 与当前的扩域 \(L\)。例如设 \(t\) 在 \(\mathbb Q\) 上超越，则 \(\mathbb Q\) 在 \(\mathbb Q(t)\) 中相对代数闭：其中在 \(\mathbb Q\) 上代数的有理函数只有有理常数，\cref{b1:ex:14-10}给出了这一事实的核验。但 \(\mathbb Q\) 本身不是代数闭域，因为 \(x^2-2\) 没有有理根。前者只问所给扩域是否带来新的代数元，后者要求 \(k[x]\) 中每个非常数多项式都有 \(k\) 中的根。

\ProseParagraphBreak
若整环 \(A\) 包含 \(k\)，且 \(A=k[a_1,\ldots,a_n]\)，其分式域是函数域。反过来，若 \(L=k(a_1,\ldots,a_n)\)，先取 \(L\) 中的子环 \(A=k[a_1,\ldots,a_n]\)。它位于域 \(L\) 中，因而是整环；由\cref{b1:prop:fraction-in-ambient-field}，\(\operatorname{Frac}(A)\) 正是 \(L\) 中包含 \(A\) 的最小子域，即
\[
L=k(a_1,\ldots,a_n)=\operatorname{Frac}\bigl(k[a_1,\ldots,a_n]\bigr).
\]
写 \(k[a_1,\ldots,a_n]\) 时，元素由生成元的 \(k\) 系数多项式给出；写 \(k(a_1,\ldots,a_n)\) 时，则还允许非零分母。二者未必相等：若 \(t\) 超越，\(1/t\in k(t)\)，却不在 \(k[t]\) 中，因为多项式 \(tq(t)\) 的常数项为零，不可能等于 \(1\)。

\ProseParagraphBreak
例如，设 \(k\) 为域、\(t\) 在 \(k\) 上超越，则 \(k[t^2,t^3]\) 的分式域仍为 \(k(t)\)，因为 \(t=t^3/t^2\)。所以不同的环可能有同一个函数域。还可考虑
\[
L=k(t)(u),\qquad u^2=t^3-t.
\]
多项式 \(x^2-(t^3-t)\) 在 \(k(t)[x]\) 中不可约。为证明这一点，假设 \(t^3-t=(A/B)^2\)，其中 \(A,B\in k[t]\) 均非零，并已将分式约成 \(\gcd(A,B)=1\) 的形式。清除分母后有
\[
A^2=(t^3-t)B^2=t(t-1)(t+1)B^2.
\]
在 \(k[t]\) 的唯一分解中，\(t\) 在 \(t^3-t\) 中恰出现一次。因此右边的因子 \(t\) 出现次数为一加上它在 \(B\) 中出现次数的两倍，是奇数；左边的出现次数则是它在 \(A\) 中出现次数的两倍，是偶数，矛盾。所以 \(t^3-t\) 不是 \(k(t)\) 中的平方。二次多项式无根便不可约，故可用不可约多项式的商域构造 \(L\)，并有 \(\bigl[L:k(t)\bigr]=2\)。

\ProseParagraphBreak
这给出一个带有两个生成元 \(t,u\)、一个关系的函数域，但不能将超越次数机械地算作“生成元数减方程数”。方程可以重复或由其他方程推出，生成元也可以冗余；独立参数的数目须由代数独立性判断。这里 \(t\) 超越，而 \(u\) 在 \(k(t)\) 上代数，才是超越次数为一的理由。刚才算出的次数二是相对于参数域 \(k(t)\) 的次数，并未判断是否存在另一参数 \(s\) 使 \(L=k(s)\)。

\subsection{次数与超越次数}
\label{b1:subsec:1404:02}

设 \(t\) 在 \(k\) 上超越，\(n\geq1\)。则 \(t^n\) 仍超越，否则将零化 \(t^n\) 的非零多项式中的变量换成 \(x^n\)，会得到零化 \(t\) 的非零多项式。并且
\begin{equation}\label{b1:eq:rational-power-degree}
\bigl[k(t):k(t^n)\bigr]=n.
\end{equation}
事实上，\(t\) 被 \(x^n-t^n\in k(t^n)[x]\) 零化，所以次数至多为 \(n\)。若 \(1,t,\ldots,t^{n-1}\) 存在非平凡的 \(k(t^n)\) 线性关系，清除分母得到
\[
\sum_{i=0}^{n-1}p_i(t^n)t^i=0,\qquad p_i\in k[x].
\]
清除分母后至少有一个 \(p_i\ne0\)。写 \(p_i(X)=\sum_jc_{ij}X^j\)，则上式变为
\[
\sum_{i=0}^{n-1}\sum_j c_{ij}t^{nj+i}=0.
\]
第 \(i\) 组中的幂指数模 \(n\) 余 \(i\)；在 \(0\le i<n\) 的范围内，不同的 \((i,j)\) 给出不同的指数 \(nj+i\)，所以不同组没有同次幂，不能相互抵消。\(t\) 超越使每个系数 \(c_{ij}=0\)，进而所有 \(p_i=0\)，与关系非平凡矛盾，证明了线性无关性。次数因而恰为 \(n\)，而 \(x^n-t^n\) 是 \(n\) 次首一零化多项式；由\cref{b1:prop:minimal-polynomial-degree}，它正是 \(t\) 在 \(k(t^n)\) 上的最小多项式，在 \(k(t^n)[x]\) 中不可约。

\ProseParagraphBreak
固定同一个纯超越扩张 \(L=k(t)/k\)。由上面的证明，\(\{t\}\)、\(\{t^2\}\) 和 \(\{t^n\}\) 都是它的超越基：这些元素在 \(k\) 上超越，且 \(t\) 在各自生成的参数域上代数。采用超越基 \(\{t\}\) 时，参数域为 \(k(t)=L\)，剩余次数 \([L:k(t)]=1\)；采用 \(\{t^2\}\) 时，参数域为 \(k(t^2)\)，剩余次数 \([L:k(t^2)]=2\)；一般地，对 \(n\ge1\)，采用 \(\{t^n\}\) 时，参数域为 \(k(t^n)\)，剩余次数 \([L:k(t^n)]=n\)。

\ProseParagraphBreak
扩张 \(L/k\) 的超越次数始终为一，剩余有限次数却随参数域而变。特别地，\([L:k(t^2)]=2\) 不能用来断定 \(L/k\) 不是纯超越扩张；\(L=k(t)\) 已经给出了反例。

\ProseParagraphBreak
当 \(\operatorname{char}k=p\) 且 \(n=p\) 时，次数仍为 \(p\)，\(x^p-t^p\) 在 \(k(t^p)[x]\) 中仍是不可约的最小多项式。把系数域扩大到 \(k(t)\) 后，才有分解 \(x^p-t^p=(x-t)^p\)，其中 \(t\) 已经属于新的系数域。不可约性必须注明系数域；这个分解也说明它只有一个不同的根，根的个数与扩张次数须加以区分，\chapref{b1:ch:16}将解释其差别。

\ProseParagraphBreak
上一小节的 \(L=k(t)(u)\) 同样满足 \(\operatorname{trdeg}_k L=1\)，因为 \(L/k(t)\) 代数，而 \(\{t\}\) 仍独立，故 \(\{t\}\) 是 \(L/k\) 的超越基。关系 \(u^2=t^3-t\) 表明 \(u\) 在这个参数域上代数，所以两个生成元只有一个独立参数。

\subsection{维数理论的域论背景}
\label{b1:subsec:1404:03}

把 \(t_1,\ldots,t_r\) 当作自由参数时，多项式环 \(k[t_1,\ldots,t_r]\) 没有参数之间的多项式关系，其分式域的超越次数就是 \(r\)。现在取 \(k[X_1,\ldots,X_n]\) 的素理想 \(\mathfrak p\)，令
\[
A=k[X_1,\ldots,X_n]/\mathfrak p.
\]
由\cref{b1:prop:prime-max-quotient}，素理想条件保证 \(A\) 是整环，因而可以取分式域。非零常数是多项式环的单位，不能属于真理想 \(\mathfrak p\)，所以 \(k\) 自然嵌入 \(A\)。剩余类 \(x_i\) 生成函数域 \(\operatorname{Frac}(A)=k(x_1,\ldots,x_n)\)。从这组生成元中抽出的超越基刻画其中相互独立的参数；其余生成元在这些参数上代数。

\ProseParagraphBreak
例如 \(k[X,Y]/(Y-X^2)\cong k[t]\)，同构由 \(X\mapsto t,Y\mapsto t^2\) 给出：用 \(Y=X^2\) 可把每个剩余类化为 \(X\) 的多项式，而这样的多项式映到零只能本来为零。商环中 \(Y\) 的剩余类已由 \(X\) 的剩余类平方得到，因此 \(Y\) 是这组生成元中的冗余项，可以消去。所得函数域 \(k(t)\) 只有一个独立参数。一般的方程也可能冗余，商环可能不再是整环，所以增加一个方程后，独立参数的变化要根据具体理想判断。

\ProseParagraphBreak
第三卷第5章将建立整扩张与 Noether 正规化，第9章将定义 Krull 维数\footnote{克鲁尔（Wolfgang Krull，1899-1971），德国数学家，发展交换环与理想理论，定义了以其名字命名的维数。}，并在域上有限生成整环的情形证明维数等于函数域的超越次数。这里尚未定义那种环的维数，因而只把超越次数视为已建立的域论不变量。
